% from ega4-2-fr.tex
\oldpage[IV-2]{36}
\section{Cycles premiers associés et décompositions primaires}

Nous donnons surtout dans ce paragraphe la traduction des résultats sur les
modules exposés dans Bourbaki, \emph{Alg. comm.}, chap.~IV, que nous suivons de
très près. Les notions qui suivent ne semblent guère avoir d'intérêt que dans
le cas de préschémas localement noethériens.

\subsection{Cycles premiers associés à un Module}

\begin{definition}[3.1.1]
\label{IV.2.3.1.1-def-fr}
Soient $X$ un préschéma, $\mathcal F$ un $\mathcal O_X$-Module
quasi-cohérent. On dit qu'un point $x\in X$ est associé à $\mathcal F$ si
l'idéal maximal $\mathfrak m_x$ de $\mathcal O_x$ est associé au
$\mathcal O_x$-Module $\mathcal F_x$ (autrement dit, si $\mathfrak m_x$ est
l'annulateur d'un élément de $\mathcal F_x$). On dénote par
$\operatorname{Ass}(\mathcal F)$ l'ensemble des $x\in X$ associés à
$\mathcal F$. On dit qu'une partie fermée irréductible $Z$ de $X$ est un
\emph{cycle premier associé à $\mathcal F$} si son point générique est
associé à $\mathcal F$. Lorsque $\mathcal F=\mathcal O_X$, on dira aussi que
les points (resp. \emph{cycles premiers}) associés à $\mathcal F$ sont
associés au préschéma $X$.

On dit qu'un cycle premier associé à $\mathcal F$ (resp. $X$) est
\emph{immergé} s'il est contenu dans un \emph{autre} cycle premier associé à
$\mathcal F$ (resp. $X$) (autrement dit, s'il n'est pas \emph{maximal} dans
l'ensemble de ces cycles).
\end{definition}

Si $X$ est localement noethérien, les \emph{composantes irréductibles} de $X$
ne sont autres que les cycles premiers \emph{maximaux} (ou \emph{non immergés})
associés à $X$.

Il est clair que si $x\in\operatorname{Ass}(\mathcal F)$, on a
$\mathcal F_x\ne0$, autrement dit

\begin{env}[3.1.1.1]
\label{IV.2.3.1.1.1-fr}
\[
\operatorname{Ass}(\mathcal F)\subset\operatorname{Supp}(\mathcal F).
\]
\end{env}

Si $x\in X$ est associé à $\mathcal F$, il est évidemment associé à
$\mathcal F|U$ pour tout voisinage ouvert $U$ de $x$, et réciproquement, s'il
est associé à $\mathcal F|U$ pour un de ces voisinages, il est associé à
$\mathcal F$.

On notera enfin que pour un $\mathcal O_X$-Module quasi-cohérent $\mathcal F$,
l'existence de cycles premiers associés immergés est une question \emph{locale},
car si $y$ et $z$ sont deux points de $\operatorname{Ass}(\mathcal F)$ tels
que $y\in\overline{\{z\}}$, tout voisinage de $y$ contient $z$.

\begin{proposition}[3.1.2]
\label{IV.2.3.1.2-fr}
Soient $A$ un anneau noethérien, $M$ un $A$-module,
$X=\operatorname{Spec}(A)$, $\mathcal F=\widetilde M$~; pour qu'un point
$x\in X$ soit associé à $\mathcal F$, il faut et il suffit que l'idéal
premier $\mathfrak j_x$ de $A$ soit associé au module $M$ (autrement dit,
soit l'annulateur d'un élément $f\in M$).
\end{proposition}

Cela résulte de la définition (3.1.1) et de Bourbaki, \emph{loc. cit.},
\S~1, n\textsuperscript{o}~2, cor. de la prop.~5, appliqué à
$S=A-\mathfrak j_x$.

\begin{proposition}[3.1.3]
\label{IV.2.3.1.3-fr}
Soient $X$ un préschéma localement noethérien, $\mathcal F$ un
$\mathcal O_X$-Module quasi-cohérent, $x$ un point de $X$, $Z$ le
sous-préschéma fermé réduit de $X$ ayant $\overline{\{x\}}$ pour espace
sous-jacent (I, 5.2.1). Les conditions suivantes sont équivalentes~:
\begin{enumerate}
\item[\emph{a)}] $x\in\operatorname{Ass}(\mathcal F)$.

